% =============================================================
% Section: Main results
% =============================================================

\section{The productivity arithmetic of alternative hours caps}\label{sec:results}

The main lesson is that a 40-hour cap and a 36-hour cap are different economic objects. Moving from 44 to 40 hours cuts the long tail of the contracted-hours distribution and lets many workers move toward the assumed efficiency peak. Moving below 40 hours is no longer mainly about trimming fatigue. It asks the economy to give up hours that the model no longer treats as inefficient. That is why the curve in Figure~\ref{fig:areq_curve} bends upward rather than sliding down smoothly.

Table~\ref{tab:main_results} reports the 36-hour endpoint, but the useful comparison is the slope of the path. The 40-hour alternative needs only about a $2\%$ one-off productivity gain to restore output. The 36-hour endpoint needs a much larger productivity response: $6.63\%$ under the one-sided fatigue benchmark and $8.18\%$ under the two-sided penalty benchmark. The distinction between these two columns is not a technicality. In the first, shorter workweeks are not assigned a direct efficiency penalty once hours fall below 40. In the second, moving away from 40 in either direction is costly.

Informality moves in the expected direction, but it is not the villain of the aggregate result. The calibrated economy does shift some marginal jobs toward informality, and the increase matters for incidence. Still, most of the output loss comes from the simpler arithmetic of fewer formal hours. The joint $(\sigma,\omega,\eta_I)$ envelope reinforces that point: changing formal-informal substitutability, the formal CES weight, and informal relative efficiency moves the endpoint requirement, but it does not turn the 36-hour reform into a small productivity problem. Capital adjustment can shave roughly one to two percentage points from the short-run benchmark over a three-to-five-year horizon, but it does not erase the difference between 40 and 36 hours.

\begin{table}[htbp]
\centering
\caption{Main results at the 36-hour endpoint.}\label{tab:main_results}
\small
\begin{tabular}{@{}l c c@{}}
\hline\hline
 & One-sided fatigue & Two-sided penalty \\
 & (preferred) & (conservative) \\
\hline
\multicolumn{3}{l}{\textit{Positive predictions (no TFP compensation)}} \\
Output, $\Delta Y$                       & $-6.60\%$        & $-8.02\%$ \\
Welfare (GHH), $\Delta$CV                & $-1.76\%$        & $-3.63\%$ \\
Informality, $\Delta$Inf                 & $+1.57$ pp       & $+1.92$ pp \\
Output per hour, $\Delta Y/h$            & $+2.39\%$        & $+0.75\%$ \\
\hline
\multicolumn{3}{l}{\textit{Benchmark: TFP gain that restores output}} \\
$A_{\text{req}}$, central                & $6.63\%$         & $8.18\%$ \\
$A_{\text{req}}$, joint envelope         & $[5.62,\,7.48]\%$ & $[6.93,\,9.23]\%$ \\
\hline\hline
\end{tabular}
\\[2pt]
\begin{minipage}{0.95\textwidth}
\justifying
{\footnotesize The joint envelope spans the eight corners of the disciplined $(\sigma,\omega,\eta_I)$ box. The one-sided fatigue envelope lies below the two-sided penalty envelope because the former does not penalize efficiency below 40 hours.}
\end{minipage}
\end{table}

\begin{figure}[htbp]
\centering
\includegraphics[width=0.92\textwidth]{fig_areq_vs_hours.pdf}
\caption{Required TFP gain $A_{\text{req}}$ as a function of the weekly hours cap.}\label{fig:areq_curve}
\vspace{2pt}
\begin{minipage}{0.92\textwidth}
\justifying
{\footnotesize \textit{Notes:} The vertical markers distinguish the 40h alternative in the relator proposal from the 36h endpoint in the original PEC 8/2025 proposal. The one-sided fatigue schedule yields $6.63\%$ at the $36$h cap. The two-sided penalty schedule yields $8.18\%$. The two schedules coincide above $h^{*}=40$ and diverge below it, reflecting different assumptions about $e(h)$ below the efficiency peak.}
\end{minipage}
\end{figure}

The figure also clarifies what the model is, and is not, saying. The first four hours of reduction are cheap in productivity units because the hours-efficiency schedule is doing helpful work. The next four hours are expensive because the reform is now mostly removing effective labor. Appendix Table~\ref{tab:efficiency_diagnostics} reports the diagnostic checks behind this statement. With no hours-efficiency channel at all, the 36-hour endpoint would require $7.74\%$ rather than $6.63\%$ or $8.18\%$; shifting the efficiency peak by one hour changes the endpoint more than changing several secondary parameters. This is why I treat the 41--42h welfare maximum below as a model result, not as a free-standing policy rule.

The historical TFP comparison should be read in the same spirit. PWT 11.0 is not a validation exercise; it is a ruler. Brazil's post-1990 TFP record is weak, so a one-off gain of six to eight percent is a large number on that ruler. The 40-hour alternative lives on a different part of the scale. Appendix Table~\ref{tab:tfp_history} translates the one-off requirements into annualized gains over three-, five-, and ten-year phase-ins. That phase-in framing is more useful than asking whether a single historical decade mechanically ``matches'' the counterfactual.

\begin{figure}[htbp]
\centering
\includegraphics[width=0.95\textwidth]{fig_areq_vs_tfp_history.pdf}
\caption{Brazilian TFP and the required one-off productivity gain.}\label{fig:tfp_history}
\vspace{2pt}
\begin{minipage}{0.95\textwidth}
\justifying
{\footnotesize \textit{Notes:} Panel~A plots the PWT 11.0 \texttt{rtfpna} series (real TFP at constant national prices) for Brazil, 1954 to 2023. Panel~B shows annualized TFP growth rates by horizon. Labels indicate the years required to accumulate the preferred 36-hour baseline $A_{\text{req}}=6.63\%$ at each historical rate (``never'' when the rate is non-positive). \textit{Source:} Penn World Table 11.0 via FRED series \texttt{RTFPNABRA632NRUG} \citep{FeenstraInklaarTimmer2015,PWT2025,FREDPWTBrazil2026}.}
\end{minipage}
\end{figure}

A downside stress test in the appendix asks what happens if the reorganization itself temporarily lowers productivity. This is not the central prediction, but it is a useful guardrail: a one-percent post-reform TFP disruption pushes the two-sided endpoint requirement above nine percent. The reason to report this case is not to forecast disruption. It is to show that the policy arithmetic is sensitive to implementation quality, not only to steady-state labor supply.
